


\section{Emerging Applications}
Verbal RL has driven improvements across multiple domains, including coding assistants, scientific discovery, and robotics, where all three pillars increasingly operate in concert. Due to page limit, we maintain focus on methodology and defer detailed discussion of applications to Appendix~\ref{app:applications} 
% Verbal RL has shown improvement in multiple domains, where all three pillars operate in concert. Due to page limit, We  page lkeep the focus of this paper on the methodology and move the application to the appendix. we detail these applications in Appendix~\ref{app:applications}. 
% Just as unsupervised learning expanded machine learning by removing dependence on labeled data, verbal reinforcement learning expands agent capabilities by making natural language the primary feedback channel. Agentic systems increasingly use outputs from tools, environments, and humans as verbal feedback to iteratively refine behavior. While the application space is rapidly growing, we highlight three domains where verbal RL has shown the strongest impact. 
% \das{shouldnt improving the model's alignment itself like training is also one of the most important application of verbal rl?}

% \subsection{Coding Assistants}
% % \das{please use newer papers, 2023 shinn is just reflexion. i dont think its the correct thing to use, you could use either things like a multi agentic blog post from anthroopic or a papers on coding agents or surveys on coding agents etc. Look at T1 paper for more references. considering how things have evolved with coding agents, 2023 is jurassic age.}
% Coding assistants are among the earliest and most successful applications of verbal RL because software environments provide deterministic feedback through compiler errors, stack traces, test failures, and linter warnings \citep{yang2024swe, chen2024teaching}. Agents translate natural-language goals into code, execute it, and iteratively refine outputs using execution feedback \citep{yang2023intercode}. Modern systems autonomously handle debugging, test generation, repository navigation, and multi-step issue resolution \citep{shinn2023reflexion}. Tool outputs such as search and dependency traces enable reasoning across large codebases, while persistent memory captures user preferences and repository constraints \citep{wang2023voyager}. SWE-bench \citep{jimenez2024swebenchlanguagemodelsresolve} has emerged as the standard benchmark for evaluating these agents.

% \subsection{Scientific Discovery and Research Workflows}
% Scientific discovery is naturally suited to verbal RL because the scientific process itself i.e. hypothesizing, experimenting, interpreting, and revising is language-driven. In drug discovery, agents iteratively refine molecular candidates using verbal feedback from chemistry and simulation tools \citep{m2024augmenting}. In materials science, agents explore compositional spaces through simulator-guided refinement \citep{szymanski2023autonomous}. Beyond individual tasks, automated systems now support end-to-end workflows including ideation, literature review, experiment design, execution, and manuscript generation \citep{lu2024ai}. Human-in-the-loop frameworks further incorporate researcher feedback during iterative refinement \citep{schmidgall2025agent}. Overall, verbal feedback enables scientific agents to continuously evaluate and revise hypotheses, accelerating research workflows. 

% \subsection{Robotics}
% Verbal RL is also transforming robotics by connecting natural-language instructions with physical control. SayCan \citep{ahn2022can} translates high-level instructions into grounded robot actions constrained by physical capabilities, while Monologue \citep{huang2022inner} incorporates environmental feedback for replanning and recovery during execution. Code as Policies \citep{liang2022code} converts language instructions into executable control code and improves behavior through execution feedback. More recent systems such as ELLMER \citep{mon2025embodied} demonstrate long-horizon task completion in dynamic environments through continual adaptation to verbal and environmental feedback.

% Industry systems are beginning to deploy these capabilities in real-world settings. NVIDIA’s GR00T N1 \citep{bjorck2025gr00t} provides a foundation model for language-conditioned robot control, while 1X Technologies’ NEO humanoid learns tasks through verbal interaction and human guidance.\footnote{\url{https://www.1x.tech}} 
% Together, these systems reflect a broader shift toward robotic agents that rely on verbal feedback to operate in open-ended environments.
% % \das{I dont have any idea what this is, and without any white paper or technical paper you cant claim how a closed source project learns. I suggest removing 1x.tech, you can maybe add openvla/octo, nvidia cosmos etc. }